Autonomous AI agents deployed in Security Operations Centres (SOCs) are
susceptible to failures arising from mismatches between task demands and
agent capabilities across skill, knowledge, and capability dimensions.
Despite growing deployment, no formal reliability model that quantifies
this mismatch and predicts failure risk before execution has, to the best
of our knowledge, been proposed. We present ARES (Agentic Reliability
Evaluation for Security systems), a formal reliability framework that
introduces three components: a five-dimensional
Skill--Knowledge--Capability (SKC) mismatch vector; a learnable Agent
Reliability Score (ARS) with task-class-specific weights; and an
Epistemic Gap metric that formalises the hallucination precondition as
confident ignorance. Experimental validation on 1,000 agent execution
traces demonstrates that knowledge mismatch is the dominant predictor of
hallucination (bucket-aggregated correlation 0.92 at $p < 0.0001$,
point-biserial correlation 0.34 at the record level), the ARS threshold
generalises across five threat classes (held-out test F1 of 0.96), and
the Epistemic Gap provides an interpretable decomposition of
false-negative risk into knowledge and calibration components.
The framework provides an initial formal foundation for pre-deployment
reliability screening and runtime anomaly detection in agentic AI systems.
